\section{Hoja de ruta de la evolución de la AGI}

\subsection{Simular el mundo $\rightarrow$ Explorar el mundo $\rightarrow$ Inducir el mundo}

Jiang Daxin (蒋大昕) propone una ruta de evolución de la AGI en tres etapas:

\begin{knowledgebox}{Marco de tres etapas}
\begin{enumerate}
\item \textbf{Simular el mundo}: representado por GPT-4O — fusión multimodal, para modelar el mundo físico
\item \textbf{Explorar el mundo}: representado por FSD V12 — control de extremo a extremo, del mundo digital al físico
\item \textbf{Inducir el mundo}: representado por O1 — pensamiento lento de Sistema 2, descubrimiento de patrones y conocimiento
\end{enumerate}
En los últimos 18 meses se lograron avances emblemáticos en las tres etapas.
\end{knowledgebox}

\subsection{Perspectiva histórica del Scaling}

Yang Zhilin (杨植麟) extrae una conclusión central de setenta u ochenta años de historia de la IA:

\begin{importantbox}{Lo único que ha funcionado en la historia de la IA es el Scaling}
Ya sea la cantidad de parámetros, la cantidad de datos o la cantidad de cómputo, más grande siempre es mejor. O1 no es solo un cambio cuantitativo, sino cualitativo: abrió una nueva dimensión de Scaling (RL Scaling) y disipó la inquietud de «qué hacer cuando se agoten los datos». Más importante aún, los datos generados por Self-play en teoría no tienen límite superior.
\end{importantbox}

\subsection{Resumen de la sección}

La evolución de la AGI puede entenderse a partir de las tres etapas «simular-explorar-inducir». El Scaling sigue siendo el motor central; O1 abrió la nueva dimensión de RL Scaling y rompió el límite del muro de datos.

